\documentclass[a4paper]{arrowhead}

\usepackage[yyyymmdd]{datetime}
\usepackage{etoolbox}
\usepackage[utf8]{inputenc}
\usepackage[T1]{fontenc}
\usepackage{multirow}
\usepackage{longtable}
\usepackage{tikz}
\usetikzlibrary{positioning,shapes.multipart,arrows.meta}

\renewcommand{\dateseparator}{-}
\setlength{\emergencystretch}{3em}
\lstset{language=HTTP,breakatwhitespace=false,columns=fullflexible}

%% Special references
\newcommand{\mref}[2][]{{\textcolor{ArrowheadPurple}{\hyperref[sec:model:#2]{\ifstrempty{#1}{#2}{#1}}}}}
\newcommand{\pref}[1]{{\textcolor{ArrowheadGrey}{\hyperref[sec:model:primitives:#1]{#1}}}}
%%

\usepackage{xurl}
\begin{document}
\sloppy

%% Arrowhead Document Properties
\ArrowheadTitle{GateTunnel}
\ArrowheadType{System Design Description}
\ArrowheadTypeShort{SysDD}
\ArrowheadVersion{5.3.0}
\ArrowheadDate{2026-11-09}
\ArrowheadAuthor{Jerker Delsing}
\ArrowheadStatus{REVIEW DRAFT}
\ArrowheadContact{jerker.delsing@ltu.se}
\ArrowheadFooter{\href{www.arrowhead.eu}{www.arrowhead.eu}}
\ArrowheadSetup
%%

%% Front Page
\begin{center}
  \vspace*{1cm}
  \huge{\arrowtitle}

  \vspace*{0.2cm}
  \LARGE{\arrowtype}
  \vspace*{1cm}

  \vspace*{\fill}

  % Front Page Image

  \vspace*{1cm}
  \vspace*{\fill}

  % Front Page Abstract
  \begin{abstract}
    This document describes the design and implementation of the GateTunnel support system of Arrowhead (documentation edition 5.3.0) (\texttt{ah5-GateTunnel-java-spring}).
  \end{abstract}

  \vspace*{1cm}
\end{center}

\newpage
%%

%% Table of Contents
\tableofcontents
\newpage
%%

\section{Documentation edition and implementation scope}
This is the 5.3.0 documentation edition, dated 2026-11-09. It carries forward the system and service contracts from the retained 5.2.2 documentation and identifies the current primary implementation sources. The former \path|documentation/v5.2.2| folder is retained unchanged.

Document version 5.3.0 is distinct from executable, dependency, registered system/service and wire-contract versions. Values still defined by the source (including 5.2.2 artefact versions, 1.0.0 service versions and \path|SysMonitor-AMQP-IDD-5.2.2| payload identifiers) are preserved. This documentation update does not change runtime behaviour or establish a new implementation release. Historical verification statements describe the tests recorded in the inherited baseline; this revision validates document compilation and references.

\section{Overview}
\label{sec:overview}

This document describes the GateTunnel system implementation, which realizes the GateTunnel SysD \cite{gt-sysd}. GateTunnel carries the inter-cloud service calls negotiated by the Gatekeepers: a consumer session listens locally and forwards each request over the relay to the provider cloud's GateTunnel, which calls the provider.
In Section \ref{Implementation}, we describe implementation details of the system.

\newpage

\section{Implementation}
\label{Implementation}

This implementation realizes the GateTunnel SysD in \texttt{ah5-\allowbreak{}GateTunnel-\allowbreak{}java-\allowbreak{}spring}. It follows the Arrowhead Java-Spring system pattern of the support systems (as Blacklist and EventHandler). The Java sources, the private copy of the Arrowhead common library (\texttt{vendor/\allowbreak{}}), tests, images and deployment definition are owned by this folder. The earlier native implementation (JDK \texttt{HttpsServer}, SQLite, package \texttt{eu.\allowbreak{}arrowhead.\allowbreak{}port}) has been replaced; the wire contracts are unchanged. \texttt{legacy/\allowbreak{}} preserves an older implementation and is excluded from the build.

\subsection{Implementation language and tools}
\begin{itemize}
\item Java 21, Spring Boot 3.4 (web, security, data JPA, Quartz), the vendored Arrowhead common library (common-utils, data-transfer-objects), Bouncy Castle, Eclipse Paho MQTT v3 (relay client) and the RabbitMQ AMQP client (sysMonitor). Maven builds the Spring Boot JAR \texttt{arrowhead-\allowbreak{}gatetunnel-\allowbreak{}5.2.2.jar}; the POM inherits only the Spring Boot parent.
\item Storage: a private MySQL 8.4 schema \texttt{ah\_\allowbreak{}gatetunnel} through JPA (MariaDB driver, \texttt{sslMode=verify-full}, \texttt{ddl-auto=validate}).
\end{itemize}

\subsection{Functional properties implementation}

The package \texttt{eu.\allowbreak{}arrowhead.\allowbreak{}gatetunnel} holds the following components:

\noindent\begin{longtable}{| p{5.6cm} | p{9.6cm} |} \hline
\rowcolor{gray!33} Component & Responsibility \\ \hline \endhead
\texttt{GateTunnel\allowbreak{}SystemInfo}, \texttt{init/} & Registered system and services; at start, the system certificate must be \texttt{GateTunnel.<cloud>.<organisation>.arrowhead.eu} with profile \texttt{sy}. Registration itself is done by the common \texttt{ApplicationInitListener}. \\ \hline
\texttt{api/\allowbreak{}http/} & Controllers of gateTunnelConnection, gateTunnelKey, gateTunnelManagement, generalManagement, echo and liveness. Each path is served under \texttt{/gatetunnel}, \texttt{/gateway} and without base path. \\ \hline
\texttt{security/\allowbreak{}GateTunnel\allowbreak{}Caller\allowbreak{}Filter} & Caller rules after the common authentication filter: certificate profile \texttt{sy}/\texttt{op}; management paths for LocalCloudManager and Sysop; other operations for Gatekeeper, DynamicServiceOrchestration, SimpleStoreServiceOrchestration, LocalCloudManager and Sysop; inter-cloud sessions for the local Gatekeeper only. \\ \hline
\texttt{service/\allowbreak{}InterCloud\allowbreak{}Session\allowbreak{}Service} & Inter-cloud relay sessions (in memory): provider side executes relayed requests within the negotiated operations; consumer side runs a mutual-TLS tunnel listener per session in the pool 8200--8299 and forwards over the relay. Ported verbatim in behaviour from 5.2.2. \\ \hline
\texttt{service/\allowbreak{}Session\allowbreak{}Security}, \texttt{PeerNames}, \texttt{validation/\allowbreak{}Operation\allowbreak{}Scope} & AH-GT-1 key exchange, sealing and replay window; explicit cross-organisation bindings; scope of a relayed request. Byte-for-byte compatible with the Go and Rust GateTunnels. \\ \hline
\texttt{service/\allowbreak{}Relay\allowbreak{}Client} & Paho MQTT client with the GateTunnel's own relay identity (\texttt{RELAY\_\allowbreak{}*}); topics \texttt{arrowhead/\allowbreak{}relay/\allowbreak{}<cloud>/\allowbreak{}gatetunnel/\allowbreak{}in/\allowbreak{}<sender CN>/\allowbreak{}<id>}. \\ \hline
\texttt{service/\allowbreak{}Provider\allowbreak{}Client} & Calls the negotiated provider with the own identity token, pinned to the provider's certificate CN, renewing a refused token once. \\ \hline
\texttt{service/\allowbreak{}Control\allowbreak{}Session\allowbreak{}Service}, \texttt{jpa/} & 4.x-compatible control sessions (control port pool 8000--8099). \\ \hline
\texttt{service/\allowbreak{}GateTunnel\allowbreak{}Jobs} & Scheduled relay (re)connection (10 s) and session expiry (60 s). \\ \hline
\end{longtable}

\subsubsection{Data handled and stored}

Storage: MySQL schema \texttt{ah\_\allowbreak{}gatetunnel}:

\noindent\begin{longtable}{| p{6.4cm} | p{8.8cm} |} \hline
\rowcolor{gray!33} Tables & Content \\ \hline \endhead
\texttt{session} & Control sessions: session and queue identifiers, peer name, service definition, reserved port, state, parties as JSON. Closed sessions are removed after \texttt{SESSION\_\allowbreak{}MAX\_\allowbreak{}AGE\_\allowbreak{}SECONDS}, which frees their ports. \\ \hline
\texttt{logs} & Log entries (Log4j2 JDBC appender), read by generalManagement \texttt{get-log}. \\ \hline
\end{longtable}

Inter-cloud sessions and their keys are kept in memory only. The interfaces exchange JSON objects; their structure is specified in the IDD documents. Error responses use the Arrowhead error body with the 5.2.2 status codes.

\subsection{Non functional properties implementation}

\subsubsection{Security}

The web server requires client certificates (PKCS12 key and trust stores) and checks them against the cloud CA's CRL. The common \texttt{OutsourcedFilter} verifies the identity token with Authentication and binds it to the client certificate; the caller filter then applies the rules above. Tunnel listeners and provider calls use the same identity and CRL outside the web server. \texttt{SecurityProfileGuard} refuses a configuration below the secure profile. The database account has data-manipulation privileges only.

\subsubsection{Power management}

None; the system is a server process without power states. The \texttt{sysMonitor} \texttt{power} observation always reports line supply.

\subsubsection{Internal monitoring}

Logs go to stdout, a log file and the \texttt{logs} table. The common \texttt{SysMonitorPublisher} publishes the \texttt{sysMonitor} observations; the health contribution is \texttt{degraded} while the relay link is down.

\subsubsection{Configuration}

\texttt{config/\allowbreak{}application.\allowbreak{}properties}; the 5.2.2 environment names (\texttt{LOCAL\_\allowbreak{}CLOUD}, \texttt{ADVERTISED\_\allowbreak{}ADDRESS}, \texttt{NEIGHBOR\_\allowbreak{}CLOUDS}, \texttt{FEDERATION\_\allowbreak{}*}, \texttt{RELAY\_\allowbreak{}*}, port pools, session times) are placeholders in it. Secrets are Compose secrets read through \texttt{configtree:/\allowbreak{}run/\allowbreak{}secrets/}.

\subsection{Deployment}

The local \texttt{Dockerfile} builds this folder and runs \texttt{arrowhead-\allowbreak{}gatetunnel-\allowbreak{}standalone:5.2.2} as UID 10001 with a read-only root filesystem. \texttt{compose.yaml} adds \texttt{gatetunnel-\allowbreak{}db} (MySQL 8.4.11, TLS required, internal network; \texttt{deployment/\allowbreak{}database/\allowbreak{}init.sql} and \texttt{least-\allowbreak{}privilege.sh}) and joins an operator-selected cloud network. The API uses port 8453; the health check is a TCP probe of it (the former plain-HTTP health port 8481 is gone). Certificates and credentials are mounted individually. Existing test-cloud deployments are not changed by this packaging.

ServiceRegistry is the bootstrap endpoint; Authentication and the orchestrator are found through it. Federation requires the organisation MQTT relay, configured neighbours, the relay identity and organisation trust. The Compose default is AH-GT-1 with explicit legacy-mTLS compatibility; peers must agree, and there is never a silent downgrade.

\subsubsection{Ownership and build boundary}
The Arrowhead common library is a private vendored source copy (\texttt{vendor/\allowbreak{}PROVENANCE.json}); security fixes must be propagated to every vendored copy. There is no shared runtime process, database or Maven parent besides Spring Boot's. Build with Java 21 and \texttt{mvn -\allowbreak{}B clean verify} (\texttt{scripts/\allowbreak{}verify.sh} also checks the JAR contents), or with the local Docker build.

\subsection{Local modifications}
\label{sec:local}

Notes on this implementation in the workspace:
\begin{itemize}
  \item Java-Spring port of the native implementation; inter-cloud sessions for Dynamic\allowbreak{}Service\allowbreak{}Orchestration (local extension).
  \item Services registered with version \texttt{5.\allowbreak{}2.\allowbreak{}2}; generalManagement adds \texttt{get-log} and \texttt{get-config}.
  \item \texttt{/health} now requires an identity token; unprefixed path aliases are served.
\end{itemize}

\subsection{Verification}

JUnit tests: the AH-GT-1 reference vector (key exchange, key derivation, sealing, replay), peer names, operation scope, TLS pinning and token renewal, the refused insecure start, and a MockMvc contract test of every served path and method through the real authentication filter. Integration tests against containers: the application against a fresh database container (schema validated with DML-only grants), and an AH-GT-1 round trip between two relay identities over an MQTT broker with mutual TLS. A live federation run with the Go and Rust GateTunnels has not been repeated for the Spring build.

\newpage

\section{Services}
\label{sec:services}

The implementation of the services is based on the SD and IDD documents referenced below.

\noindent\begin{longtable}{| p{5.4cm} | p{2.2cm} | p{2.2cm} | p{4.8cm} |} \hline
\rowcolor{gray!33}\textbf{Services produced} & \textbf{SysD ref} & \textbf{SD ref} & \textbf{IDD ref} \\ \hline \endhead
\texttt{gate\allowbreak{}Tunnel\allowbreak{}Connection} & \cite{gt-sysd} & \cite{gt-gateTunnelConnection-sd} & \cite{gt-gateTunnelConnection-http-idd} \\ \hline
\texttt{gateTunnelKey} & \cite{gt-sysd} & \cite{gt-gateTunnelKey-sd} & \cite{gt-gateTunnelKey-http-idd} \\ \hline
\texttt{gate\allowbreak{}Tunnel\allowbreak{}Management} & \cite{gt-sysd} & \cite{gt-gateTunnelManagement-sd} & \cite{gt-gateTunnelManagement-http-idd} \\ \hline
\texttt{sysMonitor} & \cite{gt-sysd} & \cite{gt-sysMonitor-sd} & \cite{gt-sysMonitor-amqp-idd} \\ \hline
&&& \\ \hline
\rowcolor{gray!33}\textbf{Services consumed} & \textbf{SysD ref} & \textbf{SD ref} & \textbf{IDD ref} \\ \hline
serviceDiscovery & \cite{core-sr-sysd} & \cite{core-sr-serviceDiscovery-sd} & -- \\ \hline
identity & \cite{core-auth-sysd} & \cite{core-auth-identity-sd} & -- \\ \hline
inter-cloud messages & Relay (MQTT) & -- & -- \\ \hline
negotiated service & Providers & -- & -- \\ \hline
sysMonitor (publisher) & AMQP broker (SysMonitor) & -- & -- \\ \hline
\end{longtable}


\section{Go and Rust implementation variants}
\subsection{Go implementation}
Primary system folder: \path|Eclipse-Arrowhead-go/ah5-GateTunnel-go|.\par
The system owns its Go module, launcher and build configuration. Use the per-system \path|Makefile| and \path|go.mod|; the containing Go tree is the build source, not a historical Arrowhead copy. Ordinary native roles use independent identities and JSON snapshot state rather than the Java database. Do not assume Java SQL migration, clustering or identical transaction/retry semantics. The registry has its own configuration and optional native AMQP monitor; other roles must be checked against their own registrations.
\paragraph{Listening port.} Built-in default 127.0.0.1:8453 (Arrowhead default). The shipped config.json, compose.yaml and Dockerfile use the same port on all interfaces ("listen": "0.0.0.0:8453"; 8540--8555 layout before 2026-11-09). Override with "listen" in the configuration file, SERVER\_PORT=N or -server.port N; the port overrides keep the configured host.
\paragraph{Native behavior and outstanding work.} Legacy-mTLS relay forwarding, owner-constrained session control, expired-session rejection. AH-GT-1 data plane; control-session rows are not additional AMQP tunnels; listeners are volatile and renegotiated on restart. These are the boundaries recorded by the native compatibility document; later implementation changes still require direct qualification.
Native design/route evidence: \path|Eclipse-Arrowhead-go/ah5-GateTunnel-go/documentation/GateTunnel|.\par
For this Java interface, select only a native transport and operation explicitly implemented and advertised by that port. A matching method/path does not establish DTO, error, authorization, retry or persistence equivalence. Java MQTT/AMQP bindings are not automatically provided by a Go or Rust role. Differential contract and security checks remain necessary before substitution.
\subsection{Rust implementation}
Primary system folder: \path|Eclipse-Arrowhead-rust/ah5-GateTunnel-rust|.\par
The per-system entry point is declared as a binary of \path|arrowhead-support-rust| in \path|ah5-support-runtime-rust/Cargo.toml|. Build from the primary Rust workspace with the applicable package/binary selection. Documentation links can resolve into the shared core/support catalogue. Independent process identity and state remain required.\par
The ordinary native profile uses Tokio/Axum and rustls with peer-identity checks and independent locked JSON state. Network effects are not distributed database transactions. The precise trust, CRL reload and persistence guarantees are those of the native profile, not the Java implementation.\par
\paragraph{Listening port.} Default 127.0.0.1:8453 (Arrowhead default, role gatetunnel). Override with "listen": "host:port" in the configuration, SERVER\_PORT=N or --server.port N; the port overrides keep the configured host.
\paragraph{Native behavior and outstanding work.} Signed AH-GT-1 exchange, dynamic mTLS listeners, encrypted forwarding, bounded sessions, replay and scoped operation guards. Exhaustive DTO/chaos tests, independent review; segment-exact grants are stricter than reference descendant grants. These are the boundaries recorded by the native compatibility document; later implementation changes still require direct qualification.
Native design/route evidence: \path|Eclipse-Arrowhead-rust/ah5-support-runtime-rust/documentation/GateTunnel|.\par
For this Java interface, select only a native transport and operation explicitly implemented and advertised by that port. A matching method/path does not establish DTO, error, authorization, retry or persistence equivalence. Java MQTT/AMQP bindings are not automatically provided by a Go or Rust role. Differential contract and security checks remain necessary before substitution.

\bibliographystyle{IEEEtran}
\bibliography{bibliography}

\newpage

\section{Revision History}
\subsection{Amendments}

\noindent\begin{tabularx}{\textwidth}{| p{1cm} | p{3cm} | p{2cm} | X | p{4cm} |} \hline
\rowcolor{gray!33} No. & Date & Version & Subject of Amendments & Author \\ \hline

1 & 2026-10-22 & 5.2.2 & Document created from the Arrowhead LaTeX templates for the 5.2.2 Java Spring implementation, including the local extensions of this workspace. & Jerker Delsing \\ \hline
2 & 2026-10-25 & 5.2.2 & Independent Java build, private helpers and state, per-system deployment and configuration. & Codex (draft) \\ \hline
3 & 2026-10-30 & 5.2.2 & Checked against the primary Java source tree (Eclipse-Arrowhead-java): standalone deployment, configuration and security statements updated. & Claude Opus 5.5 (draft) \\ \hline
4 & 2026-11-08 & 5.2.2 & Implementation sections updated for the Java-Spring port (common library, MySQL via JPA, Spring beans for the relay and AH-GT-1); contracts unchanged. & Claude Opus 5.5 (draft) \\ \hline
5 & 2026-11-09 & 5.3.0 & New documentation edition; primary-source references and Go/Rust implementation scope; 5.2.2 retained. & Documentation maintenance \\ \hline
\end{tabularx}

\subsection{Quality Assurance}

\noindent\begin{tabularx}{\textwidth}{| p{1cm} | p{3cm} | p{2cm} | X |} \hline
\rowcolor{gray!33} No. & Date & Version & Approved by \\ \hline

1 & & \arrowversion  &  \\ \hline

\end{tabularx}

\end{document}
